\chapter{From Completion to Reasoning (Statistical Manipulation)}

An LLM is a statistical engine. It outputs a probability distribution over its vocabulary (e.g., 100,277 tokens) indicating the likelihood of the next token. 

To build an agent, we must manipulate this distribution to suppress creative hallucination and force deterministic, logical output. 

\section{The Mathematics of Sampling}

When the final LayerNorm and Linear Projection are applied in the Transformer, we get raw unnormalized scores called \textbf{Logits} ($z$).

To pick the next word, we apply a Temperature-scaled Softmax:
\begin{equation}
P(x_i) = \frac{\exp(z_i / T)}{\sum_j \exp(z_j / T)}
\end{equation}

\subsection{Temperature ($T$)}
\begin{itemize}
    \item \textbf{$T = 1.0$:} Standard softmax.
    \item \textbf{$T \to 0$:} The largest logit dominates completely. The output becomes perfectly deterministic (Greedy Search).
    \item \textbf{$T > 1.0$:} The distribution flattens, making lower-probability tokens more likely.
\end{itemize}

\textbf{Agentic Implication:} For coding agents and tool-calling, \textbf{Temperature must be 0}. You want the exact function name, not a ``creative'' alternative.

\section{Logprob Entropy and Dynamic Temperature Scaling}

During long reasoning loops, setting $T = 0$ is too rigid for multi-step algorithmic planning, while setting $T = 0.7$ results in syntactically broken code. Modern runtimes solve this via \textbf{Dynamic Temperature Scaling} based on \textbf{Logprob Entropy}.

\subsection{Entropy of the Token Distribution}
The token distribution entropy $H(P)$ measures the uncertainty of the model's logits at step $t$:
\begin{equation}
H(P) = -\sum_{i \in \text{Vocab}} P(x_i) \log P(x_i)
\end{equation}

\begin{itemize}
    \item \textbf{Low Entropy ($H \to 0$):} The model is highly confident (e.g., predicting the closing bracket \texttt{\}} in JSON or completing \texttt{self.}). The runtime forces $T \to 0$ to ensure syntax validation.
    \item \textbf{High Entropy ($H > 2.0$):} The model is choosing between multiple paths (e.g., choosing which database indexing algorithm to implement). The runtime dynamically increases $T$ to allow alternative logic routes.
\end{itemize}

\begin{figure}[h]
\centering
\begin{tikzpicture}[
    scale=0.85, transform shape,
    node distance=1cm and 0.5cm,
    box/.style={draw, rectangle, rounded corners, minimum height=0.7cm, align=center, fill=blue!10},
    arrow/.style={thick,->,>=stealth}
]
\node (logits) [box] {Raw Logits};
\node (ent) [box, right=of logits] {Compute Entropy $H(P)$};
\node (scale) [box, right=of ent] {Scale Temperature $T = f(H)$};
\node (soft) [box, below=of scale] {Apply Dynamic Softmax};
\node (sample) [box, left=of soft] {Sample Token};

\draw [arrow] (logits) -- (ent);
\draw [arrow] (ent) -- (scale);
\draw [arrow] (scale) -- (soft);
\draw [arrow] (soft) -- (sample);
\end{tikzpicture}
\caption{Dynamic Temperature Sampling Flow}
\end{figure}

\subsection{Python Logprob Entropy Calculator}
Here is how the runtime evaluates entropy and scales temperature dynamically before token generation.


\begin{center}
\noindent\rule{0.6\textwidth}{0.4pt} \\
\vspace{0.3em}
\textit{[Code implementation is available in the companion coding handbook: \texttt{coding-handbook/ch02\_reasoning/}]} \\
\noindent\rule{0.6\textwidth}{0.4pt}
\end{center}


\section{Logit Bias and Schema Enforcement}

When you ask an Agent to ``output valid JSON,'' how does the API actually guarantee it? It manipulates the logits directly at the inference level.

\subsection{Structured Output (JSON Mode)}
APIs like OpenAI's Structured Outputs build a Finite State Machine (FSM) based on your JSON Schema. During generation, if the FSM determines that the next character \textit{must} be a quotation mark \texttt{"}, the engine applies a massive negative Logit Bias (e.g., $-100$) to every single token \textit{except} \texttt{"}. 


\begin{center}
\noindent\rule{0.6\textwidth}{0.4pt} \\
\vspace{0.3em}
\textit{[Code implementation is available in the companion coding handbook: \texttt{coding-handbook/ch02\_reasoning/}]} \\
\noindent\rule{0.6\textwidth}{0.4pt}
\end{center}


\section{The Chain of Thought (CoT) Exploit}

If we force the temperature to 0 and constrain the logits to JSON, the model is highly constrained. However, it still suffers from the lack of a "scratchpad". Because the model is auto-regressive, it cannot look ahead.

\subsection{Mathematical Probability Shift of CoT}
In auto-regressive generation without Chain of Thought, the model computes the probability of outputting target tokens $Y = (y_1, \dots, y_n)$ directly given prompt $X$:
\begin{equation}
P(Y | X) = \prod_{i=1}^n P(y_i | y_{<i}, X)
\end{equation}
Forcing Chain of Thought introduces intermediate reasoning tokens $Z = (z_1, \dots, z_m)$ before $Y$. The joint probability becomes:
\begin{equation}
P(Z, Y | X) = \left( \prod_{j=1}^m P(z_j | z_{<j}, X) \right) \left( \prod_{i=1}^n P(y_i | y_{<i}, Z, X) \right)
\end{equation}
The generation of $Y$ is now conditioned on $Z$. The self-attention mechanism maps $Y$'s queries against $Z$'s keys and values in the KV cache, shifting the probability distribution towards logical sanity. By the time the model generates the target response, its queries are attending to the highly logical steps inside the scratchpad block.

\section{Local LLM Grammar Enforcement via Context-Free Grammars}

While cloud providers enforce structured JSON outputs via proprietary logit biases, local deployments (using frameworks like Outlines or llama.cpp) enforce structure mathematically via \textbf{Context-Free Grammars (CFGs)}.

\subsection{FSM Logit Masking}
The runtime compiles a target schema (e.g., a Pydantic model) into a \textbf{Finite State Machine (FSM)}. At each generation step:
\begin{enumerate}
    \item The FSM inspects the current partial text buffer and identifies the state.
    \item It queries the vocabulary index to find which token strings are valid next steps (e.g., if inside a string literal, accept alphanumeric characters; if after a colon, accept a quote or bracket).
    \item It creates a binary logit mask. Disallowed tokens have their logits overwritten to $-\infty$, ensuring they have exactly $0\%$ probability of being selected by the sampler.
\end{enumerate}

\begin{center}
\noindent\rule{0.6\textwidth}{0.4pt} \\
\vspace{0.3em}
\textit{[Code implementation is available in the companion coding handbook: \texttt{coding-handbook/ch02\_reasoning/local\_grammar\_enforcement.py}]} \\
\noindent\rule{0.6\textwidth}{0.4pt}
\end{center}

\section{Reasoning Model Integration and the Planning Shift}

The introduction of native reasoning models (OpenAI o1/o3, DeepSeek-R1) changes agent design. These models run multi-step internal Chain-of-Thought planning loops *prior* to returning response tokens.

\subsection{Thoughtless Agent Loops}
Traditional ReAct orchestrators prompt the LLM to structure outputs as:
\begin{quote}
\texttt{Thought: I must call tool X... \\ Action: call\_tool(X)}
\end{quote}
With reasoning models, prompting for "Thought" blocks is redundant and degrades execution speed. Runtimes transition to \textbf{Thoughtless Loops}: the agent orchestrator simply sends the raw system prompt, and the model uses its hidden reasoning phase to solve constraints. The output contains only the final tool invocation or answer.

\subsection{XML Reasoning Tag Extraction}
Models like DeepSeek-R1 emit their reasoning process explicitly inside `<think>...</think>` XML tags. Agent runtimes must parse these tags to separate the internal reasoning (useful for system tracing and observability) from the final functional response.

\begin{center}
\noindent\rule{0.6\textwidth}{0.4pt} \\
\vspace{0.3em}
\textit{[Code implementation is available in the companion coding handbook: \texttt{coding-handbook/ch03\_react/reasoning\_model\_integration.py}]} \\
\noindent\rule{0.6\textwidth}{0.4pt}
\end{center}

